\section{Complete-source approximation at time zero}\label{sec:canonical}

We now supply all hypotheses of Theorem~\ref{thm:log} for the canonical
field at a declared height. The off-axis center supplies a safe anchor;
the observation tests still reach the real axis.

\subsection{A zero-free normalization and a safe anchor}

For $X\ge10$ set
\begin{equation}\label{eq:canonical-F}
 \begin{gathered}
 c=X+4\ii,\qquad s(z)=\frac12-\frac{\ii z}2,\\
 F(z)=\frac{s(z)-1}{s(z)}\zeta(s(z))
     =D(z)H_0(z),\qquad
 D(z)=\frac{16\pi^{s(z)/2}}{s(z)^2\Gamma(s(z)/2)} .
 \end{gathered}
\end{equation}
Throughout $\overline D(c,R)$,
\begin{equation}\label{eq:s-range}
 \frac14\le\Re s\le\frac{19}4,\qquad
 |\Re z|\le X+\frac92 .
\end{equation}
The multiplier $D$ is holomorphic and zero-free on a neighborhood
of this disk. The zeta pole at $s=1$ is cancelled in $F$, while
$s=0$ and the zeros of $1/\Gamma(s/2)$ are outside.

For $\sigma=\Re s>0$, partial summation, initially at $\sigma>1$,
and then analytic continuation give
\begin{equation}\label{eq:fractional-part}
 \frac{s-1}{s}\zeta(s)
 =1-(s-1)\int_1^\infty\{v\}v^{-s-1}\dd v .
\end{equation}
Indeed $\zeta(s)=s\int_1^\infty\lfloor v\rfloor v^{-s-1}\dd v$,
and $\lfloor v\rfloor=v-\{v\}$. The last integral in
\eqref{eq:fractional-part} converges absolutely for $\sigma>0$.
On the disk, $|s-1|\le X/2+6$, so
\begin{equation}\label{eq:canonical-upper}
 |F|\le1+\frac{|s-1|}{\sigma}\le2X+25=M_X-1,
 \qquad M_X=2X+26 .
\end{equation}
At $c$, $\Re s=5/2$. The absolutely convergent Euler product
\cite{DLMFEulerProductOC} and $|(s-1)/s|\ge3/5$ imply
\begin{equation}\label{eq:c0}
 |F(c)|\ge c_0,\qquad
 c_0=\frac35\,\frac{\zeta(5)}{\zeta(5/2)}>0 .
\end{equation}
In fact
$|\zeta(s)|\ge\prod_p(1+p^{-5/2})^{-1}
=\zeta(5)/\zeta(5/2)$, and the rational factor follows by comparing
$(3/2)^2+(\Im s)^2$ with $(5/2)^2+(\Im s)^2$.
This anchor uses an Euler product only in its absolute convergence
half-plane.

If $\widehat F$ is holomorphic near the disk and
\begin{equation}\label{eq:delta-requirement}
 \norm{F-\widehat F}_{L^\infty(D(c,R))}\le\delta
 \le\min(1,c_0/2),
\end{equation}
then $f=F/M_X$ and $g=\widehat F/M_X$ satisfy the preceding theorem
with the common constants
\begin{equation}\label{eq:canonical-constants}
 \begin{gathered}
 m_X=\frac{c_0}{2M_X},\quad A_X=\log\frac{2M_X}{c_0},\quad
 n_X=\max\!\left(1,\left\lceil\frac{A_X}{\log(27/26)}\right\rceil\right),\\
 \alpha_X=\frac2{n_X+2},\qquad C_X=C(A_X,n_X).
 \end{gathered}
\end{equation}
Consequently
\begin{equation}\label{eq:EX}
 \int_{D(c,r)}|\log|F|-\log|\widehat F||\dd A
 \le E_X,\qquad E_X=C_X(\delta/M_X)^{\alpha_X}.
\end{equation}
The factors $M_X$ cancel in the logarithm difference.
Here $n_X=O(\log X)$, $C_X=O((\log X)^2)$ and
$\alpha_X\asymp1/\log X$. For example
$A_X/n_X\le\log(27/26)$ makes the exponential in $C_X$
uniformly bounded. No height-independent positive exponent is claimed.

\subsection{Both source tails and the enlarged-disk multiplier}

For integers $J\ge1$ and $U>0$, let $\Phi_J$ be the sum in
\eqref{eq:source} restricted to $\ell\le J$, and define the exact
finite-source field
\begin{equation}\label{eq:finite-source}
 \widehat H_{J,U}(z)=\int_0^U\Phi_J(u)\cos(zu)\dd u,\qquad
 \widehat F=D\widehat H_{J,U}.
\end{equation}
These are exact finite sums and compact integrals; no quadrature or
rounding has been substituted. The real cosine source ensures
conjugation symmetry of $\widehat H_{J,U}$.

Put $b_* =17/2$ and
\begin{align}
 C_2(J)&=\sum_{\ell>J}\ell^4\ee^{-\pi(\ell^2-1)},
 & B_0&=5C_2(0),\label{eq:source-tail-defs}\\
 I_{b_*}&=\int_0^\infty
       \ee^{(b_*+1)u}r_u^2\ee^{-r_u}\dd u,
 &I_{b_*}^{>U}&=\int_U^\infty
       \ee^{(b_*+1)u}r_u^2\ee^{-r_u}\dd u .\notag
\end{align}
On the enlarged disk, $|\Im z|\le b_*$ and
$|\cos(zu)|\le\ee^{b_*u}$.
Splitting the omitted arithmetic terms on $[0,U]$ from the
\emph{complete} radial tail on $(U,\infty)$, and applying
\eqref{eq:source-basic-majorant}, proves
\begin{equation}\label{eq:delta-H}
 \sup_{|\Im z|\le b_*}|H_0(z)-\widehat H_{J,U}(z)|
 \le\delta_H(J,U):=5C_2(J)I_{b_*}+B_0I_{b_*}^{>U}.
\end{equation}
The integral is a half-line cosine integral, so there is no extra
factor of two.

Both tails have explicit majorants. For $\ell\ge J+1$, consecutive
terms in $C_2(J)$ have ratio at most
$16\ee^{-\pi(2J+3)}<1$. Therefore
\begin{equation}\label{eq:arithmetic-tail}
 C_2(J)\le
 \frac{(J+1)^4\ee^{-\pi((J+1)^2-1)}}
      {1-16\ee^{-\pi(2J+3)}}\qquad(J\ge0).
\end{equation}
The substitution $v=\pi\ee^{4u}$, with $\dd u=\dd v/(4v)$,
gives
\begin{equation}\label{eq:radial-tail}
 I_{b_*}^{>U}=\frac{\pi^{-19/8}}4
       \Gamma(35/8,\pi\ee^{4U}),\qquad
 I_{b_*}=\frac{\pi^{-19/8}}4\Gamma(35/8,\pi),
\end{equation}
where $\Gamma(a,t)=\int_t^\infty v^{a-1}\ee^{-v}\dd v$.

An explicit bound for $D$ on the whole enlarged disk is
\begin{equation}\label{eq:multiplier-bound}
 |D(z)|\le\mathcal B_X,\qquad
 \mathcal B_X=
 32\pi^{19/8}
 \exp\!\left(\frac{19}8+
       \frac{\pi^2(361+4V_X^2)}{768}\right),
 \qquad V_X=X+\frac92 .
\end{equation}
To verify it, the reciprocal-gamma product
\cite{DLMFReciprocalGammaOC} gives
\[
 \frac1{\Gamma(w)}
   =w\ee^{\gamma_{\!E}w}\prod_{k\ge1}(1+w/k)\ee^{-w/k},
 \qquad
 |\Gamma(w)^{-1}|
 \le |w|\exp\!\left(\gamma_{\!E}\Re w+
                      \frac{\pi^2|w|^2}{12}\right).
\]
Here $\log|1+q|\le\Re q+|q|^2/2$ follows from
$\log(1+t)\le t$ applied to $|1+q|^2$, and
$\sum k^{-2}=\pi^2/6$. Apply the estimate to $w=s/2$,
use $\gamma_{\!E}\le1$, $|s|\ge1/4$ and
$|s|^2\le361/16+V_X^2/4$. The prefactor in $D$ is then
$8/|s|\le32$, proving \eqref{eq:multiplier-bound}.
This bound is expensive in $X$; it is used for completeness, not
as a practical-cost assertion.

The input error for \eqref{eq:delta-requirement} is consequently
\begin{equation}\label{eq:complete-delta}
 \delta=\mathcal B_X
       \left[5C_2(J)I_{b_*}+B_0I_{b_*}^{>U}\right].
\end{equation}
Both independent tails tend to zero. Thus any positive field-error
target is attained by some finite $J,U$ at every declared $X$.
No endpoint-reflection identity for the truncated source is assumed.

